% !TEX root = ../main.tex
% =====================================================================
\section{Gleichrichter}
% =====================================================================
\begin{symtab}
U$_{di0}$ & DC-Spannung Diode & [V] \\
U$_{di\alpha}$ & DC-Spannung Tyristor & [V] \\
U$_{d}$ & DC-Spannung nach Kommutierung & [V] \\
U$_{dx}$ & DC-Spannungsabfall an induktivität & [V] \\
U$_{N}$ & Speissespannung am Gleichrichtereingang verkettet & [V] \\
P$_{N}$ & Gleichrichterleistung & [W] \\
S$_{N}$ & Gleichrichterscheinleistung & [VA] \\
i$_{d}$ & Strom durch Last & [A] \\
i$_{HL_x}$ & Strom durch einen Halbleiter & [A] \\
i$_{k}$ & Kommutierungsstrom & [A] \\
I$_{N}$ & Eingangsstrom & [A] \\
I$_{HL}$ & Halbleiterstrom & [A] \\
U$_{HL}$ & Halbleiterspannung & [V] \\
P$_{v}HL$ & Halbleiterverlustleistung (pro HL) & [W] \\
R$_{HL}$ & Halbleiterwiderstand & [$\Omega$] \\
L$_{K}$ & Leiterinduktivität & [H] \\
$\lambda$ & Leistungsfaktor & [ ] \\
$\omega$ & Kreisfrequenz & [rad$\frac{1}{s}$] \\
f$_{aus}$ & Rippelfrequenz & [Hz] \\
f$_{ein}$ & Taktfrequenz & [Hz] \\
t & Zeit & [s] \\
ü & Überlappungswinkel & [°] \\
dx & relativer induktiver Spannungsabfall & [ ] \\
$\varphi$ & Phasenverschiebungswinkel & [°] \\
\end{symtab}

\columnbreak
% ---------------------------------------------------------------------
\subsection{Einphasig (B2)}
\fig[0.6]{b2_gleichrichter}
\begin{gather*}
i_{HLrms} = \sqrt{\frac{1}{2}} \cdot i_d \\
i_{HLavg} = \frac{1}{2} \cdot i_d \\
\lambda = \frac{P_N}{S_N} = 0.9 \cdot \cos\alpha
\end{gather*}

\subsubsection{ohne Glättung}
\begin{gather*}
U_{di0} = \frac{2\sqrt{2}}{\pi} \cdot U_N = 0.9 \cdot U_N \\
U_{di\alpha} = U_{di0} \cdot \frac{1 + \cos\alpha}{2} \\
f_{aus} = 2 \cdot f_{ein}
\end{gather*}

\subsubsection{mit Glättung}
\begin{gather*}
U_{di\alpha} = U_{di0} \cdot \cos\alpha \\
P_{NX} = U_{di0} \cdot i_d = U_{di0} \cdot i_d \cdot \cos\alpha \\
S_N = U_N \cdot I_N = \frac{\pi}{2\sqrt{2}} \cdot P_N
\end{gather*}

\columnbreak
% ---------------------------------------------------------------------
\subsection{Dreiphasig (B6)}
\fig[0.6]{b6_gleichrichter}
\begin{gather*}
U_{di0} = \frac{3\sqrt{2}}{\pi} \cdot U_N = 1.35 \cdot U_N \\
U_{di\alpha} = U_{di0} \cdot \cos\alpha \\
i_N = \sqrt{\frac{2}{3}} \cdot i_d \\
P_{di\alpha} = U_{di0} \cdot i_d \cdot \cos\alpha \\
S_N = \sqrt{3} \cdot U_{Nrms} \cdot I_{Nrms} \\
= \sqrt{3} \cdot \frac{\pi}{3\sqrt{2}} \cdot U_{di0} \cdot \sqrt{\frac{2}{3}} \cdot i_d = 1.047 \cdot U_{di0} \cdot i_d \\
\lambda = \frac{P_N}{S_N} = 0.955 \cdot \cos\alpha \\
i_{HLrms} = \sqrt{\frac{1}{3}} \cdot i_d \\
i_{HLavg} = \frac{1}{3} \cdot i_d \\
P_{vHL} = U_{HL} \cdot i_{HLavg} + i_{HLrms}^2 \cdot R_{HL}
\end{gather*}

\columnbreak
% ---------------------------------------------------------------------
\subsection{Kommutierung}
\fig[0.6]{kommutierung_b6}
\fig[0.9]{kommutierung_stromform}
\begin{gather*}
\hat{I}_K = \frac{\hat{U}_N}{2 \cdot \omega \cdot L_K} \\
I_{k(t)} = \frac{\hat{U}_N}{2 \cdot \omega \cdot L_K} \cdot (1 - \cos(\omega \cdot t)) = \hat{I}_K \cdot (1 - \cos(\omega \cdot t))
\end{gather*}
Kommutierung beginnt bei $\alpha$ und endet bei $\alpha + \text{ü}$.
\begin{gather*}
\text{ü} = \arccos\left(\cos\alpha - \frac{I_d}{\hat{I}_K}\right) - \alpha \\
U_d = U_{di0}(\cos\alpha - dx) \\
U_{dx} = U_{di0} \cdot dx \\
dx = \frac{I_d}{2 \cdot \hat{I}_K}
\end{gather*}
dx normalerweise in \% angegeben.
\begin{gather*}
\varphi = \alpha + \frac{2}{3} \cdot \text{ü} \text{ bei } \alpha \lesssim 30^\circ \\
\varphi = \alpha + \frac{1}{2} \cdot \text{ü} \text{ bei } \alpha \sim 30^\circ - 60^\circ \\
\varphi = \alpha + \frac{1}{3} \cdot \text{ü} \text{ bei } \alpha \sim 60^\circ - 90^\circ
\end{gather*}

\columnbreak
% ---------------------------------------------------------------------
\subsection{Harmonische}
{\footnotesize\setlength{\tabcolsep}{3pt}
\renewcommand{\arraystretch}{1.0}
\noindent Bei 6-pulsigen Schaltungen kommen folgende Harmonische vor:\\
\begin{tabular*}{\linewidth}{@{}l@{\extracolsep{\fill}}llllllr@{}}
$U_d$ & n*6:      & 6.     & 12.      & 18.      & 24.      & 30.      & etc. \\
$i_N$ & n*6 $\pm$ 1: & 5., 7. & 11., 13. & 17., 19. & 23.,25. & 29.,31. & etc. \\
\end{tabular*}\\[4pt]
Bei 12-pulsigen Schaltungen kommen folgende Harmonische vor:\\
\begin{tabular*}{\linewidth}{@{}l@{\extracolsep{\fill}}llllllr@{}}
$U_d$ & n*12:      & & 12.      & & 24.      & & etc. \\
$i_N$ & n*12$\pm$1: & & 11., 13. & & 23.,25. & & etc. \\
\end{tabular*}\\[4pt]
Bei 24-pulsigen Schaltungen kommen folgende Harmonische vor:\\
\begin{tabular*}{\linewidth}{@{}l@{\extracolsep{\fill}}llllllr@{}}
$U_d$ & n*24:      & & & & 24.      & & etc. \\
$i_N$ & n*24$\pm$1: & & & & 23.,25. & & etc. \\
\end{tabular*}\par}
